\newpage
\section{Preliminary calculations}
The collection contains $N=1{,}000{,}000$ documents. The same 200 query vectors, 256-bit
codes and reference answers are used throughout. Each method may choose its own cluster
count $C$, probe count $P$ and local settings. A probe opens one cluster. The result count
is $K=100$; $d$ denotes the number of document bits.

\subsection{Time complexity}
We count the work for one query. Let $N$ be the number of documents, $C$ the number of
clusters, and $P$ the number opened. For this first estimate, assume equal-sized clusters:
\[
 n=N/C\quad\text{documents per cluster}.
\]
The examples below use $N=1{,}000{,}000$, $C=100$, and $P=10$. Each cluster has 10,000
rows, so the opened clusters contain 100,000 rows. These are arithmetic examples, not
additional benchmark runs. The experiment uses actual cluster sizes.

Let $T_0$ be time spent choosing clusters and preparing the query's score table.
Let $a_A,a_B,a_C$ be the average time per scored document for each method, including
reading its code and maintaining the best results. These costs are measured; scattered
reads can cost more than consecutive reads. The vector length $d$ and result count $K$
are fixed in the examples.

\subsubsection{Original method}
\textbf{Loop:} open each selected cluster, then score every document inside it.
There are $P$ clusters with approximately $N/C$ documents each:
\[
 \boxed{T_A\approx T_0+\frac{PN}{C}a_A.}
\]
For example, A scores $10\times10{,}000=100{,}000$ documents. Following Exa's
four-coordinate lookup-table approach~\cite{exa}, each 256-bit document needs 64 table
terms. Thus this example evaluates $100{,}000\times64=6{,}400{,}000$ terms, in addition
to selecting the best results. Four packed words do not mean four scoring operations.

For fixed $C$, $P$, $d$, and $K$, scanning work is $O(N)$: doubling the documents
approximately doubles the number of scores. For unequal clusters, replace $PN/C$ by
the actual count $F_A=\sum_{c\text{ opened}}n_c$.

\subsubsection{Bitplanes}
\textbf{Loop:} visit a branch, read its bitmap to split the candidates, and continue.
When a candidate set is small enough, score its remaining documents.

A bitmap stores one bit per document. With 64 bits per word, a cluster of 10,000
rows needs $\ceil{10{,}000/64}=157$ words. Every split processes all 157 positions,
even after only 50 candidates remain. Each position requires reads, bit operations
and writes; it is not a single CPU instruction.

Let $S$ be the average number of splits visited per opened cluster, including
alternative branches. Let $f_B$ be the fraction of opened documents eventually scored,
and $b$ the time to process one bitmap position during a split. Then
\[
 \boxed{T_B\approx T_0+PS\ceil{\frac{N}{64C}}b
              +f_B\frac{PN}{C}a_B+T_{\rm other,B}.}
\]
Here $T_{\rm other,B}$ is additional time for sorting query positions, managing waiting branches
and reading final masks.

For example, if $S=20$ and $f_B=0.1$, B processes
$10\times20\times157=31{,}400$ bitmap positions and scores 10,000 documents.
A scores 100,000. B wins only if its extra work costs less than the 90,000 avoided scores,
while keeping the required recall.

The bitmap term still grows with $N$ when $C$, $P$ and $S$ stay fixed. Branching therefore
does not automatically give $O(\log N)$ search. The values of $S$ and $f_B$ must be measured
on the actual queries; the example does not predict them.

\subsubsection{Backward Walk}
\textbf{Loop:} generate a key, look it up in each opened cluster, then score the documents
listed under that key. Repeat for another key if needed.

Let $E$ be the number of key patterns tried, $c$ the average lookup time, and $F_C$ the
number of documents found and scored. Trying every key in every opened cluster gives
\[
 \boxed{T_C\approx T_0+PE\,c+F_Ca_C+T_{\rm other,C}.}
\]
$T_{\rm other,C}$ is additional time for generating and ordering keys and starting document ranges.
Empty lookups also cost time. If a budget stops midway through a key, use the actual
lookup count instead of $PE$.

An $h$-bit key has $2^h$ possible patterns. \emph{If} rows are evenly distributed among
these patterns inside each cluster, $E$ distinct keys return approximately
\[
 F_C\approx\frac{PEN}{C2^h},\qquad
 \boxed{T_C\approx T_0+PE\,c+\frac{PEN}{C2^h}a_C+T_{\rm other,C}.}
\]
For example, with $h=8$ and $E=16$, C performs $10\times16=160$ lookups. Under that
assumption, it scores approximately $160\times10{,}000/256=6250$ documents.
If all 256 keys are needed, it makes 2560 lookups and scores all 100,000 documents.

Finding 6250 documents does not establish recall: the required reference documents must
be among them. Real keys can also contain very different numbers of rows. The experiment
therefore measures $E$ and $F_C$ rather than assuming this distribution.

\subsection{Memory complexity}
For cluster $c$, let $n_c$ be its document count and $W_c=\ceil{n_c/w}$ its bitmap width,
where a word holds $w=64$ bits. Stored index data and temporary query memory are separate. The running program, unused
array capacity and allocation overhead also consume RAM. Construction may have a higher
peak because input arrays and the new index coexist.

\subsubsection{Original method}
A 256-bit code occupies 32 bytes. An eight-byte ID brings the total to 40 bytes per row,
or 40 MB for one million rows. With $u=w/8$ bytes per packed word,
\[
 M_{\rm rows}=N\left(u\ceil{d/w}+b_{\rm id}\right).
\]
Cluster storage is additional. For example, $C$ floating-point centroids of width $d_R$
require $4Cd_R$ bytes for coordinates, plus labels and row ranges.

A query also needs its vector, score table, selected cluster IDs and top-$K$ results.
The current query vector occupies 1024 bytes; its 64 groups of 16 double-precision table
entries occupy 8192 bytes. These are data sizes, not the complete process footprint.

\subsubsection{Bitplanes}
B retains the rows and adds one bitplane per coordinate:
\[
 M_{\rm planes}=du\sum_{c=1}^{C}\ceil{n_c/w}.
\]
One global cluster adds 32 MB of planes to the 40 MB of rows and IDs. Separate clusters
introduce padding because each plane ends on a whole-word boundary.

Waiting branches retain masks. If $A_c(t)$ masks for cluster $c$ exist at time $t$, their
peak data size is
\[
 M_{\rm masks}=u\max_t\sum_c A_c(t)W_c.
\]
This includes a parent while both children are created. Queue entries and unused capacity
are additional. Counting only the mask currently being processed is insufficient.

\subsubsection{Backward Walk}
C stores rows consecutively within each key range. Their existing IDs serve as the document
lists, so a second $N$-element ID array is unnecessary. Each occupied entry stores a key,
a start position and a row count. For $U$ entries,
\[
 M_{\rm entries}=U(b_{\rm key}+2b_{\rm offset}).
\]
A four-byte key and two eight-byte positions require 20 bytes of field data per entry.
Hash-table pointers and allocation overhead are additional. Query memory includes the key
queue. More key bits can increase directory size without reducing the final scoring work.

\subsection{Recall estimation}
Returning 99 of the reference top 100 IDs gives 99\% recall. The reference is the exact
ranking under the unchanged full binary-vector score, with the same ID rule for ties.
This measures agreement with that ranking, not whether the text answers a person's question.

Let $G_q$ be the reference IDs and $O_{a,q}$ the IDs returned by method $a$. Batch recall is
\[
 \overline R_a=\frac{1}{n_qK}\sum_{q=1}^{n_q}|G_q\cap O_{a,q}|,
 \qquad n_q=200.
\]
This uses the saved IDs directly. It requires no assumption that document signs or embedding
coordinates are independent.

\subsubsection{Original method}
Let $D_q$ contain the documents in the opened clusters and $g_q=|G_q\cap D_q|$. A scores
all of $D_q$, so its per-query recall is exactly $R_{A,q}=g_q/K$. Opening clusters containing
95 reference IDs gives 95\% recall; scoring their other documents cannot recover the five
reference IDs in unopened clusters.

\subsubsection{Bitplanes}
Among the $g_q$ reference IDs that reached the clusters, let $r_{B,q}$ be the fraction
returned. Then $R_{B,q}=(g_q/K)r_{B,q}$. For example, if eight of ten reference documents
reach the clusters and six of those eight are returned, recall is $(8/10)(6/8)=60\%$.
Unlimited search with safe bounds makes $r_{B,q}=1$. A finite node budget must be checked
against the returned IDs. If $g_q=0$, recall is zero.

\subsubsection{Backward Walk}
Similarly, $R_{C,q}=(g_q/K)r_{C,q}$. Visiting all necessary keys, or safely ruling out the
remaining keys, makes the local fraction one. Candidate or lookup limits can reduce it.
The batch value averages these per-query products; multiplying separately averaged fractions
is generally incorrect. A key returning many rows establishes no recall level unless the
required reference IDs are among them.

\subsection{Summary}
Each formula uses its method's own opened clusters and local work. Let $M_0$ be allocated
row and router storage, $Q_0$ common query memory, $D_C$ allocated directory memory, and
$Q_B,Q_C$ extra masks and queues. Allocation overhead is included in these memory terms.

\subsubsection*{Time complexity}
\begin{center}\small
\begin{tabularx}{\linewidth}{@{}lY@{}}\toprule
Method & Approximate query time\\\midrule
Original method & $T_0+(PN/C)a_A$\\
Bitplanes & $T_0+PS\ceil{N/(64C)}b+f_B(PN/C)a_B+T_{\rm other,B}$\\
Backward Walk & $T_0+PE\,c+F_Ca_C+T_{\rm other,C}$; for equal rows per key, $F_C\approx PEN/(C2^h)$\\\bottomrule
\end{tabularx}\normalsize
\end{center}
$N$ is document count; $C$ is clusters; $P$ is opened clusters. B visits an average of $S$
splits per cluster and scores fraction $f_B$. C tries $E$ keys of width $h$.
The $a$ terms price complete document scores; $b$ prices a split-word visit; $c$ prices a
key lookup. $T_{\rm other,B}$ and $T_{\rm other,C}$ retain other search work, rather than treating it as free.
These estimates use equal cluster sizes. The measured work uses actual sizes and key counts.

$T_0$ includes routing and table preparation. Each cost belongs in one term only.
Query encoding and later reranking are outside the measured retrieval time.

\subsubsection*{Memory complexity}
\begin{center}\small
\begin{tabular}{@{}lll@{}}\toprule
Method & Stored index & Temporary query memory\\\midrule
Original method & $M_0$ & $Q_0$\\
Bitplanes & $M_0+M_{\rm planes,allocated}$ & $Q_0+Q_B$\\
Backward Walk & $M_0+D_C$ & $Q_0+Q_C$\\\bottomrule
\end{tabular}\normalsize
\end{center}
The process, including these arrays and the running program, must fit 32 GB. A peak that
includes construction is reported as such, not as index size.

\subsubsection*{Recall estimation}
\begin{center}\small
\begin{tabularx}{\linewidth}{@{}lYY@{}}\toprule
Method & Per-query recall & Local exactness condition\\\midrule
Original method & $g_q/K$ & Score all opened documents\\
Bitplanes & $(g_q/K)r_{B,q}$ & Complete search with safe bounds\\
Backward Walk & $(g_q/K)r_{C,q}$ & Visit or safely rule out remaining keys\\\bottomrule
\end{tabularx}\normalsize
\end{center}
The reported batch value averages the same 200 reference answers.
